\clearpage\pagestyle{otMain}\setlength{\otBlocEspacement}{2.5mm}\fontsize{11}{15}\selectfont\corrigesChapitre{12}{Masses et contenances}
\begin{blocLivre}{cor-12-1}\begin{corrige}{1}{La bonne unité}\demarche{Passer par des objets de référence.}
Gomme : $\res{\text g}$. Camion : $\res{\text t}$. Cuillère : $\res{\text{mL}}$. Piscine : des milliers de litres, donc $\res{\text L}$ (ou l’hectolitre).
\end{corrige}\end{blocLivre}
\begin{blocLivre}{cor-12-2}\begin{corrige}{2}{Convertir des masses}\demarche{Utiliser 1 kg = 1 000 g dans les deux sens.}
$3\times1\,000=\res{3\,000}$ g et $2{,}5\times1\,000=\res{2\,500}$ g. Pour 500 g : la moitié de $1\,000$ g, soit $\res{0{,}5}$ kg.
\end{corrige}\end{blocLivre}
\begin{blocLivre}{cor-12-3}\begin{corrige}{3}{Convertir des contenances}\demarche{Citer la relation utilisée pour chaque conversion.}
$2\times100=\res{200}$ cL. $350$ cL $=\res{3{,}50}$ L, car 100 cL font 1 L. $5\times100=\res{500}$ L.
\end{corrige}\end{blocLivre}
\begin{blocLivre}{cor-12-4}\begin{corrige}{4}{Comparer}\demarche{Convertir dans la plus petite unité.}
$1{,}2$ kg $=1\,200$ g et $1\,200>950$ : $\res{1{,}2\ \mathrm{kg}>950\ \mathrm g}$. $1$ L $=100$ cL et $75<100$ : $\res{75\ \mathrm{cL}<1\ \mathrm L}$.
\end{corrige}\end{blocLivre}
\begin{blocLivre}{cor-12-5}\begin{corrige}{5}{Des estimations sensées}\demarche{Éliminer les associations invraisemblables.}
Orange : $\res{200\ \mathrm g}$. Cartable : $\res{7\ \mathrm{kg}}$. Canette : $\res{33\ \mathrm{cL}}$. Baignoire : $\res{150\ \mathrm L}$.
\end{corrige}\end{blocLivre}
\begin{blocLivre}{cor-12-6}\begin{corrige}{6}{La balance à plateaux}\demarche{Additionner les masses marquées.}
L’équilibre signifie l’égalité des masses : $500+200+50=\res{750}$ g.
\label{fig-cor-12-6}\figureGeo[corrige]{mas-6}
\end{corrige}\end{blocLivre}
\begin{blocLivre}{cor-12-7}\begin{corrige}{7}{Le colis complet}\demarche{Convertir avant d’additionner.}
$1{,}2$ kg $=1\,200$ g ; total : $1\,200+850=\res{2\,050}$ g, soit $\res{2{,}050}$ kg, c’est-à-dire $2{,}05$ kg.
\end{corrige}\end{blocLivre}
\begin{blocLivre}{cor-12-8}\begin{corrige}{8}{Des verres bien remplis}\demarche{Chercher combien de fois 25 tient dans 150.}
$1{,}5$ L $=150$ cL. On compte de 25 en 25 : $25\times6=150$, donc $\res6$ verres exactement.
\end{corrige}\end{blocLivre}
\begin{blocLivre}{cor-12-9}\begin{corrige}{9}{Le chargement du camion}\demarche{Comparer le total à la tonne convertie.}
$850+325=1\,175$ kg, or 1 t $=1\,000$ kg. Le chargement dépasse de $1\,175-1\,000=\res{175}$ kg : $\res{\text{impossible}}$.
\end{corrige}\end{blocLivre}
\begin{blocLivre}{cor-12-10}\begin{corrige}{10}{La recette doublée}\demarche{Convertir le paquet en grammes.}
Doubler : $250\times2=\res{500}$ g. Le paquet contient $1\,000$ g : il reste $1\,000-500=\res{500}$ g, la moitié du paquet.
\end{corrige}\end{blocLivre}
\begin{blocLivre}{cor-12-11}\begin{corrige}{11}{Les briques de lait}\demarche{Calculer le stock avant la consommation.}
Le stock vaut $3\times6=18$ briques, soit 18 L. Il reste $18-15=\res3$ L, trois briques pleines.
\end{corrige}\end{blocLivre}
\begin{blocLivre}{cor-12-12}\begin{corrige}{12}{L’ascenseur}\demarche{Additionner toutes les masses embarquées.}
Les personnes pèsent $4\times72=288$ kg ; avec le colis : $288+15=\res{303}$ kg. Comme $303>300$, c’est $\res{\text{interdit}}$, à 3 kg près.
\end{corrige}\end{blocLivre}
\begin{blocLivre}{cor-12-13}\begin{corrige}{13}{La citerne de l’école}\demarche{Convertir les hectolitres avant de soustraire.}
$5$ hL $=500$ L. Il reste $500-120=\res{380}$ L.
\end{corrige}\end{blocLivre}
\begin{blocLivre}{cor-12-14}\begin{corrige}{14}{Les pots de confiture}\demarche{Chercher le facteur qui complète 300.}
$2{,}4$ kg $=2\,400$ g. En comptant de 300 en 300 : $300\times8=2\,400$, donc $\res8$ pots, sans reste.
\end{corrige}\end{blocLivre}
\begin{blocLivre}{cor-12-15}\begin{corrige}{15}{Plumes ou plomb ?}\demarche{Distinguer la masse de l’encombrement.}
Les deux pèsent $\res{\text{pareil}}$ : un kilogramme reste un kilogramme, seule la place occupée change. Un litre d’eau pèse environ $\res{1\ \mathrm{kg}}$ : cette référence relie contenances et masses pour estimer.
\end{corrige}\end{blocLivre}
